% ============================================================================
%  4. DISCUSSION
%
%  Nota de estilo de A&A: la revista publica una lista de correcciones
%  idiomáticas frecuentes. Las más habituales en manuscritos de dinámica
%  relativista:
%
%    "depend on", no "depend of"          "independent of", no "independent on"
%    "associated with", no "associated to" "comparable to", no "comparable with"
%    "typical of", no "typical for"        "adjacent to", no "adjacent with"
%    "on the order of" (aprox.), no "of order"
%    "in the past 5 years", no "in the last 5 years"
%    "such as", no "like for example"      "we describe X in detail", no "in details"
%    "in the case of", no "in case of"     "on a timescale of", no "in a timescale"
%    "originate in", no "originate from"   "to stem from", mejor que "to result from"
%    "search for X", no "search of X"      "least-square technique" (con guion)
%    "Galactic bulge", "Galaxy" en mayúscula si son la nuestra
%    Evitar "quite / rather / somewhat": el adjetivo solo suele bastar.
%    "but" y "however" dicen lo mismo: usar uno de los dos, no ambos.
%
%  Alto/bajo vs. grande/pequeño vs. fuerte/débil:
%    high/low     -> value, rate, redshift, temperature, luminosity, velocity,
%                    energy, density, eccentricity, inclination, flux
%    large/small  -> scale, amplitude, mass, momentum, uncertainties
%    short/long   -> time, length, timescale
%    strong/weak  -> correlation, gradient, dependence, constraint, instability
% ============================================================================

\section{Discussion}
\label{sec:discussion}

Texto de la sección.
